\chapter{Source Code and Reproduction}
\label{app:code}

\section{Repository}
\label{app:repo}

The complete implementation of this study is held in a version-controlled
public repository:

\begin{center}
\fbox{\parbox{0.88\textwidth}{\centering\vspace{2pt}
\textbf{Source code repository}\\[4pt]
\url{\repourl}\\[4pt]
\small Source code licensed under the \repolicence.\\
No patient data is included.
\vspace{2pt}}}
\end{center}

The Apache-2.0 grant covers the pipeline source written for this study. It does
not extend to material this work does not own, and the repository's
\texttt{NOTICE} file draws that boundary explicitly: the institutional thesis
template and its brand assets are reproduced under their owner's rights and are
excluded; the two clinical databases are excluded absolutely, being redistributed
by neither this repository nor its licence; and the thesis text and its figures,
being scholarly prose rather than software, are made available for citation and
academic reuse with attribution rather than under a software licence.

Where a routine implements a published method (the SOFA score, the Sepsis-3
operational criteria, NEWS2, qSOFA, SIRS, the PhysioNet Utility Score, the
sample-size criteria, pooled logistic regression, decision-curve analysis) the
implementation is original to this work and written from the published
definition, and the method is attributed to its authors both in \texttt{NOTICE}
and at the point of use in this thesis. No implementation was copied or adapted
from a third party's source.

The repository contains \textbf{code only}. No MIMIC-IV or eICU-CRD record,
derived person-hour table, or fitted model object is included: the
\texttt{data/}, \texttt{processed\_data/} and \texttt{output/} directories are
excluded from version control by design, because the PhysioNet Credentialed
Data Use Agreement (Section~\ref{sec:ethics}) prohibits redistribution of the
underlying data. Exclusion is enforced by \texttt{.gitignore} and was verified
against the tracked file list before release.

\section{Repository Layout}

{\small
\begin{longtable}{p{4.6cm}p{8.4cm}}
    \caption{Repository layout.}\label{tab:repo_layout}\\
    \toprule
    \textbf{Path} & \textbf{Contents} \\
    \midrule
    \endfirsthead
    \toprule
    \textbf{Path} & \textbf{Contents} \\
    \midrule
    \endhead
    \bottomrule
    \endfoot

        \texttt{00\_preregistration.md} &
        The pre-registration, including the G0 analysis gate and Protocol
        Amendments 1--8, each dated and disclosed as post-hoc \\
        \addlinespace
        \texttt{config.R} &
        All constants: label-variant parameters, feature lists, plausibility
        ranges, seeds, bootstrap sizes, multiplicity families, event floors.
        No numeric constant is defined anywhere else \\
        \addlinespace
        \texttt{utils.R} &
        Shared helpers, including \texttt{require\_features()} (errors on a
        declared-but-absent predictor) and \texttt{check\_multinom\_fit()} \\
        \addlinespace
        \texttt{clinical\_scores.R} &
        SOFA, NEWS2, qSOFA and SIRS implementations \\
        \addlinespace
        \texttt{01\_setup.Rmd} \dots\ \texttt{12\_inference.Rmd} &
        The twelve pipeline stages described in Chapter~\ref{ch:design}, in
        execution order \\
        \addlinespace
        \texttt{run\_pipeline.sh}, \texttt{run\_pipeline.R} &
        Orchestration. Runs all stages, or a named subset, and regenerates the
        thesis constants on success \\
        \addlinespace
        \texttt{thesis\_constants.R} &
        Reads the result tables and emits every numerical value used in this
        document as a \LaTeX{} macro (Section~\ref{app:repro}) \\
        \addlinespace
        \texttt{utility\_score.py} &
        Reference implementation of the Utility Score, maintained as a mirror
        of the R scoring rule and cross-checked against it \\
        \addlinespace
        \texttt{check\_consistency.R} &
        Cross-file assertions over the result tables: arithmetic identities
        that must hold between quantities written by different stages, and
        predicates for the comparative claims this document states in words \\
        \addlinespace
        \texttt{check\_thesis\_numbers.R} &
        Checks that every result-shaped number in the thesis sources is either
        generated, cited to prior work, or recorded as a design constant \\
        \addlinespace
        \texttt{check\_thesis\_margins.R} &
        Measures the built PDF for content passing the text block, which
        \LaTeX{} does not report for a table \\
        \addlinespace
        \texttt{output/} &
        Created at run time: \texttt{processed\_data/}, \texttt{figures/},
        \texttt{logs/}. Not tracked \\
\end{longtable}
}

\section{Reproducing the Results}
\label{app:repro}

\paragraph{Prerequisites.} R 4.5 or later with the packages listed in the
repository, Python 3.14 with \texttt{pandas} 3.0.4 for the Utility Score
mirror, and a \TeX{} distribution providing \texttt{xelatex} and
\texttt{pdflatex}/\texttt{biber}. The two source databases must be obtained
independently through PhysioNet credentialing: MIMIC-IV v3.1 and eICU-CRD
v2.0. They are placed under \texttt{data/} and located automatically, or
pointed at with the \texttt{MIMIC\_RESEARCH\_ROOT} environment variable.

\paragraph{Full run.} A complete execution is a single command,

\begin{center}\texttt{./run\_pipeline.sh}\end{center}

\noindent
which runs stages 01--12 in order and, on success, regenerates the thesis
constants and figures. Individual stages or phases can be run by number. Every
stochastic step is seeded, so a repeated run reproduces the reported estimates
exactly; the one exception is noted below.

\paragraph{How the thesis stays aligned with the pipeline.} No quantity
estimated by this study is typed by hand into this document; the exceptions are
figures quoted from prior work, constants fixed by the pre-registration, and
prose roundings of a macro stated exactly nearby. On completion,
\texttt{thesis\_constants.R}
reads the result tables and writes every reported quantity as a macro into
\texttt{pipeline\_constants.tex}, which the thesis includes; it also copies the
generated figures into the document's image directory. A quantity the pipeline
did not produce typesets as a conspicuous marker rather than as a stale number,
so a partial run cannot silently leave an outdated value in the text. This is
the mechanism by which the numbers in Chapters~\ref{ch:evaluation}
and~\ref{ch:discussion} are guaranteed to be the numbers the code produced.

\paragraph{Determinism.} All seeds are fixed in \texttt{config.R} and the
train/test partition is defined by admission year rather than by sampling, so
it is invariant. Gradient-boosted tree fitting is the one step that is not
bit-reproducible across differing library builds; the pipeline therefore
reports the gradient-boosted results as an explicit consistency check on any
re-run, and every other estimate is invariant to it.

\section{Ethical and Data-Access Statement}

Both databases are de-identified and publicly available to credentialed
researchers under the PhysioNet Credentialed Health Data Use Agreement. The
analysis is retrospective and observational, involves no patient contact and no
intervention, and required no additional ethical approval beyond the data use
agreement, whose terms, including the prohibition on redistribution and on
attempted re-identification, were observed throughout.
Section~\ref{sec:ethics} sets out the full statement.


\chapter{Pipeline Stage Reference}
\label{app:stages}

Chapter~\ref{ch:design} describes what each stage does and why the architecture
takes the shape it does. This appendix records the file-level interface: what
each stage reads and what it writes; the mapping from an analysis to its
inference status is in the analysis register (Section~\ref{sec:register})
rather than repeated here. Paths are relative to \texttt{sepsis/project/}; all intermediate tables
are Parquet, with a CSV mirror written for any result table small enough to read
directly.

\begin{table}[H]
    \centering\scriptsize\setlength{\tabcolsep}{4pt}
    \begin{tabularx}{\textwidth}{p{0.9cm}p{4.2cm}X}
        \toprule
        \textbf{Stage} & \textbf{Reads} & \textbf{Writes} \\
        \midrule
        01 & MIMIC-IV source tables &
        n/a (validation only; halts on a missing table) \\
        \addlinespace
        02 & \texttt{hosp/}, \texttt{icu/} source tables &
        \texttt{02\_cohort\_\{A..F\}}, \texttt{02\_label\_variance\_table},
        \texttt{02\_alt\_label\_variance\_table} \\
        \addlinespace
        03 & \texttt{02\_cohort\_*}; \texttt{chartevents}, \texttt{labevents},
        \texttt{inputevents}, \texttt{outputevents} (streamed) &
        \texttt{03\_person\_hours\_\{A..F\}}, \texttt{03\_early\_onset\_shift} \\
        \addlinespace
        04 & \texttt{03\_person\_hours\_*} (three columns) &
        \texttt{04\_sample\_size} \\
        \addlinespace
        05 & \texttt{03\_person\_hours\_*} &
        \texttt{05\_model\_primary\_*}, \texttt{05\_coef\_all\_variants},
        \texttt{05\_predictions\_primary\_*}, \texttt{05\_h2\_stability\_*};
        \texttt{*\_ablated} counterparts \\
        \addlinespace
        06 & \texttt{03\_person\_hours\_*} &
        \texttt{06\_model\_gbt\_*}, \texttt{06\_predictions\_comparators\_*},
        \texttt{06\_gbt\_importance\_*}, \texttt{06\_schoenfeld\_*};
        ablated GBT counterparts \\
        \addlinespace
        07 & \texttt{05\_*}, \texttt{06\_*} predictions &
        \texttt{07\_metric\_results}, \texttt{\_random}, \texttt{\_altlabel},
        \texttt{07\_utility\_sweep}, \texttt{07\_utility\_max},
        \texttt{07\_ablation\_metrics} \\
        \addlinespace
        08 & eICU-CRD source tables; frozen \texttt{05\_}/\texttt{06\_} models &
        \texttt{08\_external\_validation\_results}, \texttt{\_abxonly},
        \texttt{08\_eicu\_predictions\_*}, \texttt{08\_eicu\_coverage},
        \texttt{08\_event\_rate\_reconciliation},
        \texttt{08\_h5\_internal\_external} \\
        \addlinespace
        09 & \texttt{07\_metric\_results*} &
        \texttt{09\_variance\_decomposition\_table},
        \texttt{09\_hypothesis\_verdicts},
        \texttt{09\_label\_anchor\_attribution},
        \texttt{09\_label\_agreement}, \texttt{09\_pretreatment\_onsets} \\
        \addlinespace
        10 & \texttt{03\_person\_hours\_*}, \texttt{05\_}/\texttt{06\_}
        predictions &
        Per-stratum performance and label-sensitivity tables \\
        \addlinespace
        11 & \texttt{05\_}/\texttt{06\_}/\texttt{08\_} predictions &
        \texttt{11\_calibration\_summary}, \texttt{11\_operating\_points},
        \texttt{11\_net\_benefit\_at\_prevalence},
        \texttt{11\_patient\_level\_metrics},
        \texttt{11\_lead\_time\_summary}, external counterparts \\
        \addlinespace
        12 & Stored predictions only (refits nothing) &
        \texttt{12\_confirmatory\_tests}, \texttt{12\_internal\_ci},
        \texttt{12\_subgroup\_auroc\_ci}, \texttt{\_contrasts},
        \texttt{12\_subgroup\_labelsens\_*}, \texttt{12\_altlabel\_*},
        \texttt{12\_h2\_*}, \texttt{12\_analysis\_register} \\
        \bottomrule
    \end{tabularx}
    \caption{Stage inputs and outputs. Which result each stage supports is
    recorded once, in the analysis register of
    Section~\ref{sec:register}, and is not repeated here. Stage 03 is
    the only stage that reads the full-size source tables, and Stage 12 is the
    only one that reads nothing but stored predictions, which is what allows
    the inference layer to be re-run and checked without refitting any model.}
    \label{tab:stage_reference}
\end{table}

\paragraph{Naming conventions.} A file's numeric prefix is the stage that wrote
it. A \texttt{\_\{A..F\}} suffix denotes the label variant; \texttt{\_ablated}
denotes the Protocol Amendment 5 feature arm; \texttt{\_random} denotes the
random-split arm used only for H4. Files without a variant suffix are pooled
across variants by construction. No stage overwrites a file written by another
stage.
